\documentclass[12pt]{amsart}
\begin{document}

\title{Quantile tomography: using quantiles with multivariate data}
\author{Ivan Mizera}
\address{University of Alberta}
\maketitle
%\begin{abstract}
Directional quantile envelopes---essentially, depth contours---are a possible way to condense the directional quantile information, the information carried by the quantiles of projections. In typical circumstances, they allow for relatively faithful nd straightforward retrieval of the directional quantiles, offering a straightforward probabilistic interpretation in terms of the tangent mass at smooth boundary points. They can be viewed as a natural, nonparametric extension of ``multivariate quantiles'' yielded by fitted multivariate normal distribution, and, as illustrated on data examples, their construction can be adapted to elaborate frameworks---like estimation of extreme quantiles, and directional quantile regression---that require more sophisticated estimation methods than simply evaluating quantiles for empirical distributions. Their estimates are affine equivariant whenever the estimators of directional quantiles are translation and scale equivariant; mathematically, they express the dual aspect of directional quantiles.
%\end{abstract}
\end{document}

